\section{An independently checked source reduction}
\label{sec:source229-independent}

The source reduction replaces four occurrences of the register-program
subtree
\[
 A=((\mathrm{CD}_6,(\mathrm{CI}_5,\mathrm{CC}_2)),\mathrm{CI}_5)
 \qquad\hbox{by}\qquad
 B=(\mathrm{CI}_5,(\mathrm{CD}_6,\mathrm{CC}_2)).
\]
A pair denotes a cons cell in the binary source tree; $\mathrm{CI}_r$
increments register $r$, $\mathrm{CD}_r$ decrements a positive register
and breaks once when it is zero, and $\mathrm{CC}_k$ continues $k$
levels.  Register values are natural numbers.  This removes one
instruction leaf per occurrence, reducing the literal source from
5,125 leaves to 5,121.  Source leaves are not Turing-machine states;
the state saving is a separate property of the compiled table.

\subsection{An unbounded local proof}

Let $p$ be the prefix of a replaced subtree.  The old instruction
addresses relative to $p$ are
$00:\mathrm{CD}_6$, $010:\mathrm{CI}_5$,
$011:\mathrm{CC}_2$, and $1:\mathrm{CI}_5$.
The new addresses are $0:\mathrm{CI}_5$,
$10:\mathrm{CD}_6$, and $11:\mathrm{CC}_2$.
These addresses permit a direct proof in the path semantics of
\texttt{RM.v}, without appealing to an informal interpretation of
the parentheses.  Paths are read from left to right and padded by
zeros; each primitive operation manipulates the actually consumed
prefix.  Increment means fixed-length binary increment.  Continue
strips bits from the right until its specified number of ones has
been removed, while a zero decrement strips one zero and increments.

Write the input values as $R_5=s$ and $R_6=n$.  Both blocks have
exactly the following recurrence:
\begin{center}
\begin{tabular}{c|c|c|c}
Guard & Primitive steps & Resulting values & Control result \\
\hline
$n=0$ & $2$ & $(R_5,R_6)=(s+1,0)$ & common normal successor \\
$n\geq1$ & $3$ & $(R_5,R_6)=(s+1,n-1)$ & restart this block
\end{tabular}
\end{center}
Indeed, in $A$, a positive decrement is followed by an increment
and a continue to the prefix $p0$, which resolves to $p00$.
A zero decrement jumps to $p1$, increments, and returns by
$\operatorname{inc}(p1)$.
In $B$, the increment comes first.  A positive decrement is followed
by a continue to $p$, which resolves to $p0$; a zero decrement
returns by the same $\operatorname{inc}(p1)$.

Induction on $n$ therefore proves that both blocks terminate normally
after exactly $3n+2$ register-machine leaf steps, with
\[
 R_5'=R_5+R_6+1,\qquad R_6'=0,
 \qquad R_i'=R_i\quad(i\ne5,6).
\]
The induction also proves that neither block visits a halt leaf or
escapes by an unmatched break or continue.  This is an argument for
all natural register values, with no finite register bound.

\subsection{Whole-source control boundaries}

The independent script \texttt{audit\_source229.py} reconstructs both
complete binary trees from their literal path catalogues.  It verifies
that these four disjoint prefixes, and no other source positions, are
changed:
\begin{center}\small\ttfamily
\begin{tabular}{l}
111010010010111111000011\\
11101001001011111110011\\
1110100101100011\\
111110011
\end{tabular}
\end{center}
It computes the successor of every unchanged instruction directly
from the path arithmetic in \texttt{RM.v}, including both guards of
every decrement.  All 5,109 exterior leaves are identical.  All 6,635
exterior-to-exterior nonhalting edges are identical.  Exactly four
edges enter a rewritten block, each at its first instruction in
both versions.  The unique halt leaf remains at
\texttt{1110100101101}, and the initial root entry is unchanged.
These checks do not assume that a guard is reachable only under the
intended enumerator invariant.

At a common exterior instruction, equal registers therefore execute
the same operation and select corresponding successors.  At a paired
block entry, the preceding induction restores equal registers at the
same exterior successor after the same finite number of steps.
The resulting relation on exterior configurations and paired block
entries proves whole-program halting equivalence, from the root and
any natural initial register file.  Totality of the blocks ensures
that no infinite computation disappears inside an omitted block.
The compressed execution also preserves the number of register-machine
leaf steps to the halt leaf, whenever that leaf is reached.

\subsection{Connecting the JSON source to the formal theorem}

The separate parser \texttt{audit\_prog230\_bridge.py} reads the
literal \texttt{p230\_tab} in the pinned \texttt{Prog230.v}.  It checks
the exact mapping from compact edges to \texttt{RM.v} instructions,
root zero, the 5,124-node bound, the out-of-range default, and the
literal initial register expression
\texttt{upd (upd regs0 0 3) 4 3}.  It then expands the rooted table
and compares every ordered leaf against both the published
\texttt{machine/sites230.json} and the normalized source used here.

All 5,124 nodes are reachable without cycles or repeated visits;
the expansion has 5,125 leaves and maximum depth 40.  Thus the
formal program's walk fuel is sufficient for every instruction
pointer, including pointers padded by zeros.  Every expanded leaf
matches both JSON catalogues exactly.  The manifest hashes of
\texttt{Prog230.v}, \texttt{RM.v}, and the published catalogue match
the manifest retained by the successful upstream Linux proof build.
This establishes an explicit source-object correspondence, rather
than relying on a generator comment or a filename.

Together, the literal parser, the exterior graph audit, and the
unbounded block induction connect the reduced source to the actual
$p_{230}$ occurring in K230's formal theorem.  They do not by
themselves prove that a new binary Turing table implements the
reduced source: that requires the separately supplied compiler,
bootstrap, and universal runtime argument.  The present source
transport is an ordinary mathematical proof with executable finite
certificates; no new Rocq theorem for the reduced source or table
is claimed.
